% Tabel portrait, boleh lintas halaman, caption di atas header
% Perbandingan deskriptif tiga model kematangan DataOps. Tanpa skor lintas model
% karena ketiganya menilai objek berbeda (lihat prosa II.1.1)
{
\footnotesize
\setlength{\tabcolsep}{3pt}
\renewcommand{\arraystretch}{1.2}
\setlength{\LTpost}{0pt}
\Needspace{13\baselineskip}
\begin{longtable}{|>{\raggedright\arraybackslash}m{2.8cm}|>{\raggedright\arraybackslash}m{4.4cm}|>{\raggedright\arraybackslash}m{3.4cm}|>{\raggedright\arraybackslash}m{2.4cm}|}
\caption{Perbandingan Model Kematangan DataOps}
\label{tab:dataops_maturity_comparison} \\
\hline
\centering\arraybackslash\textbf{Model} & \centering\arraybackslash\textbf{Tahap atau \textit{Level}} & \centering\arraybackslash\textbf{Fokus Penilaian} & \centering\arraybackslash\textbf{Lingkup dan Jenis Sumber} \\
\hline
\endfirsthead
\hline
\centering\arraybackslash\textbf{Model} & \centering\arraybackslash\textbf{Tahap atau \textit{Level}} & \centering\arraybackslash\textbf{Fokus Penilaian} & \centering\arraybackslash\textbf{Lingkup dan Jenis Sumber} \\
\hline
\endhead
\endfoot
\hline
\endlastfoot

Model evolusi DataOps \autocite{munappy2020adhoc} & Lima tahap: \textit{ad-hoc data analysis}, \textit{semi-automated data analysis}, \textit{agile data science}, \textit{continuous testing and monitoring}, dan DataOps & Evolusi kapabilitas \textit{data pipeline}, dari pengumpulan data manual hingga orkestrasi dan pemantauan seluruh \textit{data lifecycle} & Lintas domain, studi kasus telekomunikasi, makalah telaah sejawat \\
\hline
\textit{DataOps Maturity Model} \autocite{datakitchen2020maturity} & Lima \textit{level} per dimensi: \textit{struggle}, \textit{basic}, \textit{consistent}, \textit{quantitative}, dan \textit{optimized} & Asesmen organisasi pada enam dimensi: \textit{error rates}, \textit{cycle time}, \textit{collaboration}, \textit{measurement}, \textit{team culture}, dan \textit{customer happiness} & Analitik data umum, \textit{white paper} vendor \\
\hline
Model kematangan DataOps industrial \autocite{harrington2021highbyte} & Empat tahap: \textit{data access}, \textit{data contextualization}, \textit{site visibility}, dan \textit{enterprise visibility} & Cakupan integrasi data operasional, dari akses data mesin, kontekstualisasi dan normalisasi data, hingga visibilitas lintas situs dengan \textit{data governance} & Data industrial manufaktur, artikel vendor \\
\hline
\end{longtable}
\par{\footnotesize Model-model di atas menilai objek yang berbeda sehingga pemetaan skor kapabilitas lintas model seperti pada Tabel~\ref{tab:mlops_maturity_comparison} tidak diterapkan. Perbedaan fokus tiap model dibahas pada teks.}
}
